%% ============================================================
%%  Chapter 6 – Conclusion
%% ============================================================
\chapter{Conclusion}
\label{ch:conclusion}

This thesis presented \medxplain, an end-to-end, hospital-deployable framework
for Explainable Medical Visual Question Answering. The system addresses the
dual requirements of clinical AI adoption --- privacy compliance and
transparent reasoning --- as first-class architectural concerns rather than
afterthoughts.

%% ─────────────────────────────────────────────────────────────
\section{Summary of Contributions}
\label{sec:summary}

Five original contributions were made:

\begin{enumerate}[label=\textbf{C\arabic*.}]

  \item \textbf{Hospital-Ready DICOM Ingestion Pipeline.}
        The \code{data\_ingestion/dicom\_pipeline.py} module provides
        production-grade loading of single-frame (CR, DX) and multi-frame
        (CT, MR) DICOM studies, handling all common transfer syntaxes,
        applying VOI-LUT windowing to produce normalised float32 tensors,
        and pseudonymising patient identities via deterministic SHA-256 hashing.

  \item \textbf{HIPAA/GDPR-Compliant PHI Scrubber.}
        The \code{phi\_scrubber.py} module removes all 18 HIPAA Safe Harbour
        identifiers from DICOM headers (aligned with PS~3.15 Annex~E) and
        redacts PHI from free-text radiology reports using regex and
        optional spaCy NER, maintaining a tamper-evident audit log.

  \item \textbf{Two-Phase Training Protocol.}
        Phase~1 aligns image and text embedding spaces with NT-Xent
        contrastive loss on frozen BiomedCLIP and BioGPT backbones.
        Phase~2 fine-tunes for VQA by unfreezing the last two vision encoder
        blocks and injecting LoRA adapters (rank~$r=8$) into the language
        decoder, reducing trainable parameters by 99.4\%.

  \item \textbf{Integrated Grad-CAM Explainability.}
        Heatmaps are computed and rendered for every inference within the
        Gradio clinical interface, enabling clinicians to visually verify
        the spatial evidence underlying each answer, without additional
        latency or user intervention.

  \item \textbf{Extended Clinical Evaluation Suite.}
        BERTScore F1 and per-pathology Clinical F1 (computed via the
        \code{ReportParser} pathology catalogue) join standard exact-match
        accuracy and BLEU-4, providing a richer and more clinically
        meaningful characterisation of model performance. A radiologist-comparison
        module generates PHI-free discrepancy PDFs for expert review.
\end{enumerate}

%% ─────────────────────────────────────────────────────────────
\section{Answers to Research Questions}
\label{sec:answers}

\begin{description}
  \item[\textbf{RQ1}] \textit{Can a multimodal VQA model trained on
        de-identified hospital data achieve $>$65\% accuracy?}

        Experimental results (Table~\ref{tab:main_results}) confirm / do not
        yet confirm this target. \todo{Insert actual result here.}
        The two-phase training strategy demonstrably improves over a
        single-phase baseline (Table~\ref{tab:abl_phase}).

  \item[\textbf{RQ2}] \textit{How can PHI compliance and explainability be
        first-class architectural concerns?}

        PHI scrubbing is enforced at the data ingestion layer, before any
        model sees the data, and is independently auditable. Grad-CAM
        is computed unconditionally for every forward pass and surfaced in
        the primary UI panel, not as an optional add-on.

  \item[\textbf{RQ3}] \textit{Does two-phase training outperform end-to-end
        training under a hospital GPU budget?}

        The ablation study (Table~\ref{tab:abl_phase}) shows that Phase~1
        pre-alignment consistently improves Phase~2 VQA accuracy.
        LoRA at rank~8 provides the optimal parameter-efficiency trade-off
        (Table~\ref{tab:abl_lora}).
\end{description}

%% ─────────────────────────────────────────────────────────────
\section{Future Work}
\label{sec:future}

\begin{description}
  \item[Multi-institutional federated validation.]
        Extending \medxplain{} to a federated learning setting~\citep{mcmahan2017federated}
        would allow model training across multiple hospital sites without
        centralising patient data, directly addressing the single-institution
        limitation identified in \S\ref{sec:limitations}.

  \item[Open-ended report generation.]
        Replacing the closed-answer classification head with a generative
        BeamSearch decoder would allow the model to produce full structured
        radiology reports from images, building on recent work in automated
        report generation~\citep{tanida2023rgrg}.

  \item[Prospective clinical study.]
        Measuring \emph{downstream} impact on radiologist reporting time,
        missed finding rates, and inter-reader variability in a prospective
        randomised study is the gold standard for demonstrating clinical value.

  \item[Edge / point-of-care deployment.]
        Quantising the model with INT8 post-training quantisation and
        compiling with TensorRT would enable deployment on edge GPUs in
        radiology rooms or emergency departments with minimal compute
        infrastructure.

  \item[Multi-modal fusion of structured EHR data.]
        Incorporating patient demographics, lab values, and prior diagnoses
        from the Electronic Health Record alongside the imaging data could
        further improve diagnostic accuracy, following the context-aware
        answering module already prototyped in
        \code{vqa\_app\_deliverable/advanced\_features.py}.

  \item[Longitudinal change detection.]
        The existing longitudinal view tab in the Gradio UI could be
        strengthened by a dedicated siamese-network change detection module
        that localises interval change between serial studies.
\end{description}

%% ─────────────────────────────────────────────────────────────
\section{Closing Remarks}
\label{sec:closing}

The convergence of large-scale vision-language pretraining, parameter-efficient
fine-tuning, and mature open-source tooling has made hospital-grade Medical VQA
achievable within the constraints of a single research group's GPU budget.
\medxplain{} demonstrates that privacy compliance and explainability need not
be obstacles to AI adoption in clinical settings --- they can be engineered
into the system from the ground up, making the resulting tool more trustworthy,
auditable, and deployable than post-hoc patched alternatives.

The open-source release of the full codebase at
\url{https://github.com/Darkfazer/medxplain-phase1} is intended to lower
the barrier for other research groups and hospital partners to build upon
this infrastructure for their own cohorts and clinical tasks.
